% =====================================================================================
%  Lecture 6  —  30 September 2026   (video -cqba5tE6aQ, 1 h 10 min)
% =====================================================================================
\chapter{How a language model is built and trained}\label{ch:L6}
\setlecturevideo{-cqba5tE6aQ}
\lectureinfo{30 September 2026}{1 h 10 min}{-cqba5tE6aQ}{LLM architecture, deployment and training stages}

% macros local to this chapter
\newcommand{\Lpiref}{\pi_{\mathrm{ref}}}
\newcommand{\Lrelu}{\mathrm{ReLU}}
\newcommand{\Lcal}{\mathcal{L}}

\begin{glance}
\begin{itemize}
  \item \textbf{What an LLM is.} A neural network seen as a chain of vector spaces; the \emph{transformer}: tokens $\to$ embeddings $\to$ $L$ layers of causal attention and MLP, all of one width $\to$ unembedding and \emph{softmax with temperature} $\to$ a sampled token that is appended to the text; and the loop repeats.
  \item \textbf{What really goes in; agents.} A hidden \emph{system prompt}, the conversation so far, and \emph{reasoning tokens} between a hard-coded ``think'' token and a learned end token, hidden or summarized for a safety reason and a commercial one. An agent is an LLM plus a \emph{harness}; four layers of monitoring.
  \item \textbf{Training.} Pre-training (next-token log-loss), mid-training, post-training (SFT, RLHF/RLAIF, RL with verifiable rewards); compute shifting towards RL; data and Chinchilla; why the next-token loss is \emph{not} shallow statistics (Kraft--Shannon again).
  \item \textbf{Post-training mathematics.} The SFT loss mask; $\Pr(w\succ l)=\sig\big(r(w)-r(l)\big)$; reward minus a KL penalty; DPO.
  \item \textbf{Safety.} RLHF does not scale and breeds sycophancy; Claude's Constitution versus OpenAI's Model Spec and deliberative alignment; how they are trained in; the evidence; a slide clarified.
\end{itemize}
\end{glance}

\begin{caution}[About the source for this lecture]
The recording (the session was held online) has English auto-captions only, and the slides (the speaker mentions ``slide~25'') are not visible, so every figure I describe comes from his words. The captions garble technical terms (\emph{unembedding} appears as ``unicetted'', RLHF as ``RHF'', mid-training as ``midra mid-training''), product names (Claude is often ``cloth'' or ``clot'') and several numbers; I corrected what I could identify with confidence and marked the rest\flag. Students are not named. The lecture is a descriptive tour (``feel free to correct me if I'm saying something wrong \ldots{} I am myself somewhat new to the area''); where a statement is loose, and for all checks of numbers and of real-world claims, see the ``Beyond the lecture'' boxes. In the course outline this lecture opens the week on \emph{LLM training stages}.
\end{caution}

\begin{figure}[tb]
\centering
\mapbar{4248}{600}{%
  \mapseg{0}{510}{slate}\mapseg{510}{945}{teal}\mapseg{945}{1670}{accent}\mapseg{1670}{2115}{amber}
  \mapseg{2115}{2670}{crimson}\mapseg{2670}{2870}{plum}\mapseg{2870}{3215}{moss}\mapseg{3215}{4248}{ink}
  \mapin{255}{6.1}\mapin{727}{6.2}\mapin{1307}{6.3}\mapin{1892}{6.4}\mapin{2392}{6.5}\mapin{2770}{6.6}\mapin{3042}{6.7}\mapin{3731}{6.8 \ Safety side}
  \maplabel{0}{Neural nets}{3.2}{south west}
  \maplabel{510}{Softmax, temperature}{7.2}{south west}
  \maplabel{945}{System prompt, CoT}{3.2}{south west}
  \maplabel{1670}{Agents, monitors}{7.2}{south west}
  \maplabel{2115}{Training stages}{3.2}{south west}
  \maplabel{2670}{Compression}{7.2}{south west}
  \maplabel{2870}{RLHF, DPO}{3.2}{south west}
  \maplabel{3215}{Constitution, Model Spec}{7.2}{south west}
}
\caption{\textbf{Lecture map} (1\,h\,10\,min, to scale). The first half describes the \emph{machine} (6.1--6.4: architecture, output, what goes in, agents), the second its \emph{training} (6.5--6.7) and the safety-motivated documents (6.8).}
\end{figure}

%% ------------------------------------------------------------------------------------
\section{What a language model is made of}\label{sec:L6-arch}
\tsm{0:00}
After several weeks of talking about AI and LLMs, the instructor says, it is worth asking what an LLM actually \emph{is}. The lecture goes through the basics of the architecture and of the training, partly for those who have not met them, and partly so that those who have can pick out the parts of this very large system that matter most for safety. He is himself somewhat new to the area and asks to be corrected.

\subsection{Neural networks: graphs and chains of vector spaces}\label{sec:L6-nn}\tsm{1:10}
A large language model is a neural network, and a neural network can be pictured in two ways. As a \emph{graph}: neurons are nodes joined by edges, and each edge carries a weight, the strength of the connection between two neurons. Or as a \emph{chain of vector spaces and maps}: the neurons are organized into layers, the activations $h_1,h_2,\dots$ of a layer are the components of a vector, and the connections between consecutive layers form a matrix. Passing from one layer to the next, the signal is multiplied by the matrix, a \emph{bias vector} is added, and a nonlinear function is applied to every component.

\begin{defn}[A layer of a network]\label{def:L6-layer}
With activations $h^{(\ell)}\in\R^{n_\ell}$, a weight matrix $W^{(\ell)}\in\R^{n_{\ell+1}\times n_\ell}$ and a bias vector $b^{(\ell)}\in\R^{n_{\ell+1}}$,
\[
  h^{(\ell+1)}\;=\;\varphi\big(W^{(\ell)}h^{(\ell)}+b^{(\ell)}\big),\qquad \varphi=\Lrelu,\ \ \Lrelu(s)=\max(s,0),
\]
with $\varphi$ applied to each component.
\end{defn}
\noindent The widely used nonlinearity is the \term{ReLU}: zero for negative input, the identity for positive input (the captions say ``zero at zero''). A neuron with negative input is not activated; with positive input it is activated, with a strength equal to the signal received. Some nonlinearity is essential, for otherwise the whole network would only produce linear maps from input to output, and ``the world is much more complicated than the linear maps''; this is in some ways the simplest way for a nonlinearity to enter, and it turns out to be enough to describe arbitrarily complicated functions.

A particular kind of network is the \term{multilayer perceptron} (MLP): the first model of a neuron was historically called the perceptron, and many layers of them make an MLP, also called a \term{fully connected} network because the weights connect every neuron of a layer to every neuron of the next. It is a sequence of transformations from one vector space to the next: a matrix, a pointwise nonlinearity, a matrix, a pointwise nonlinearity, and so on.

\subsection{The transformer}\label{sec:L6-transformer}\tsm{4:55}
The architecture used in LLMs is a more complicated one, the \term{transformer}: each layer contains an MLP, but also other operations. Here is how it works (Figure~\ref{fig:L6-pipeline}). The input is a sentence, or any text. It is broken into little pieces, words or parts of words, called \term{tokens}; punctuation and a line break can be tokens too. The tokens are \term{embedded}: each token is associated with a vector in the first vector space. The \term{vocabulary}, the set of all tokens, typically has around $100$ to $200$ thousand tokens, and the vector space has a dimension typically between a thousand and perhaps ten thousand, so there are one or two orders of magnitude more tokens than dimensions. The embedding starts out random and is learned, in the training, so as to be as good as possible.

The vectors then pass through the layers. The operations applied to them are MLPs and \term{attention}. Attention, the speaker says, is multiplication by matrices, ``so it is a purely linear operation'' (a loose description; see the box below), and it gives the vector of each token information about the vectors of the \emph{previous} tokens. It is \term{causal}: the information always goes from earlier tokens to later ones, and a token receives nothing from the future. This is supposed to help the model understand how the vector of a token was determined by what happened before it. The output of a layer is again a vector in the same space, so all layers have the same dimension; one can think of what happens inside an LLM as a collection of vector spaces and operations between them. Each piece is not too hard to understand, he says, even though what happens as a whole is very complicated, and he urges the class to watch the videos on attention linked from the reading material.

\begin{figure}[tb]
\centering
\begin{tikzpicture}[x=1cm, y=1cm, >={Latex[length=2mm,width=1.4mm]},
    blk/.style={card=#1, minimum height=1.55cm, text width=2.2cm, inner sep=3pt},
    lab/.style={font=\sffamily\fontsize{6.6}{7.4}\selectfont, text=slate, align=center}]
  % ---------------- row 1: from text to the block stack
  \node[blk=slate] (tok) at (2.05,0.55) {\textbf{tokens}\\[-1pt]{\scriptsize text $\mapsto t_1,\dots,t_n$\\ $V\approx10^5$--$2\!\cdot\!10^5$ tokens}};
  \node[blk=teal] (emb) at (5.05,0.55) {\textbf{embeddings}\\[-1pt]{\scriptsize $x_i^{(0)}=E[t_i]\in\R^{d}$\\ $d\approx10^3$--$10^4$}};
  \node[draw=accent, fill=accent!6, rounded corners=2pt, line width=0.7pt, minimum width=5.1cm, minimum height=3.7cm, anchor=west] (stack) at (7.1,-0.2) {};
  \node[font=\sffamily\scriptsize\bfseries, text=accent, anchor=south west] at ([xshift=2mm,yshift=1.2mm]stack.south west) {$L$ layers, one width $d$};
  \node[card=accent, minimum width=1.9cm, minimum height=0.95cm] (att) at (8.4,0.55) {causal\\attention};
  \node[circle, draw=slate, fill=white, inner sep=0.4pt, font=\sffamily\scriptsize] (p1) at (9.7,0.55) {$+$};
  \node[card=teal, minimum width=1.4cm, minimum height=0.95cm] (mlp) at (10.7,0.55) {MLP};
  \node[circle, draw=slate, fill=white, inner sep=0.4pt, font=\sffamily\scriptsize] (p2) at (11.75,0.55) {$+$};
  \draw[->, accent, line width=0.9pt] (att.east) -- (p1.west);
  \draw[->, accent, line width=0.9pt] (p1.east) -- (mlp.west);
  \draw[->, teal, line width=0.9pt] (mlp.east) -- (p2.west);
  \draw[->, slate, line width=0.55pt] (7.3,0.55) |- (9.7,1.25) -- (p1.north);
  \draw[->, slate, line width=0.55pt] (9.84,0.55) |- (11.75,1.45) -- (p2.north);
  % the causal mask icon
  \begin{scope}[shift={(7.4,-0.5)}, x=0.17cm, y=0.17cm]
    \foreach \i in {0,...,4} \foreach \j in {0,...,4} {
      \ifnum\j>\i \else \fill[accent!55] (\j,-\i) rectangle ++(0.92,-0.92); \fi
      \draw[ink!30, line width=0.2pt] (\j,-\i) rectangle ++(0.92,-0.92);}
  \end{scope}
  \node[lab, anchor=west, text width=3.5cm, align=left] at (8.45,-0.95) {mask $M$: token $i$ attends to tokens $j\le i$ only};
  \draw[->, line width=1pt, ink] (tok.east) -- (emb.west);
  \draw[->, line width=1pt, ink] (emb.east) -- (emb.east -| stack.west);
  % ---------------- row 2: from the last vector back to the text
  \node[blk=moss] (hh) at (11.35,-4.2) {\textbf{last vector}\\[-1pt]{\scriptsize $h=x^{(L)}_n\in\R^{d}$}};
  \node[blk=amber] (unemb) at (8.35,-4.2) {\textbf{unembedding}\\[-1pt]{\scriptsize $z=Uh+b\in\R^{V}$\\ rows $u_v$: one per token}};
  \node[blk=plum] (sm) at (5.35,-4.2) {\textbf{softmax}\\[-1pt]{\scriptsize $\Pr(v)\propto e^{z_v/\tau}$\\ temperature $\tau$}};
  \node[blk=crimson] (samp) at (2.35,-4.2) {\textbf{sample}\\[-1pt]{\scriptsize next token $t_{n+1}$}};
  \draw[->, line width=1pt, ink] (hh.north |- stack.south) -- (hh.north);
  \draw[->, line width=1pt, ink] (hh.west) -- (unemb.east);
  \draw[->, line width=1pt, ink] (unemb.west) -- (sm.east);
  \draw[->, line width=1pt, ink] (sm.west) -- (samp.east);
  \draw[->, line width=1.2pt, crimson] (samp.west) -- ++(-0.65,0) |- node[lab, pos=0.25, anchor=east, rotate=90, yshift=1.5mm] {append, repeat} (tok.west);
  \node[lab, text width=3.4cm] at (2.2,-1.6) {the loop of Section~\ref{sec:L6-loop}: one pass per new token};
\end{tikzpicture}
\caption{\textbf{A decoder-only transformer at a glance} (schematic; sizes as quoted in the lecture). A text is cut into tokens, each token becomes a vector of dimension $d$, and the vectors pass through $L$ layers, each a causal-attention step and an MLP step whose outputs are added to the running vector (the circled plus signs: the ``residual'' connections; see the box \emph{The transformer block in formulas}). The vector of the \emph{last} position is turned into one score per vocabulary token, then into probabilities; a token is sampled and appended, and the process repeats.}
\label{fig:L6-pipeline}
\end{figure}

\begin{beyond}[The transformer block in formulas, and what the lecture glossed over]\label{box:L6-block}
\emph{Notation} (local to this chapter): $n$ tokens, vocabulary size $V$, width $d$, $L$ layers, $H$ heads of width $d_k=d/H$. The embedding matrix $E\in\R^{V\times d}$ gives $x_i^{(0)}\in\R^d$ (plus position information, Section~\ref{sec:L6-loop}), and layer $\ell$ of a standard decoder-only transformer (in the pre-normalization form of most current models; the original \citep{vaswani2017} is an encoder--decoder with post-normalization) \emph{adds} its output to a running vector, the \term{residual stream}:
\[
  u^{(\ell)}=x^{(\ell-1)}+\mathrm{MHA}_\ell\big(\mathrm{LN}(x^{(\ell-1)})\big),\qquad
  x^{(\ell)}=u^{(\ell)}+\mathrm{MLP}_\ell\big(\mathrm{LN}(u^{(\ell)})\big).
\]
Here $\mathrm{LN}$ is a normalization and one attention head, acting on the rows of its input $X\in\R^{n\times d}$, is
\[
  \mathrm{head}(X)=\softmax\!\Big(\frac{QK^{\top}}{\sqrt{d_k}}+M\Big)\,XW_{v},\qquad Q=XW_q,\quad K=XW_k,
\]
the softmax acting on each row, with the causal mask $M_{ij}=0$ for $j\le i$ and $M_{ij}=-\infty$ for $j>i$. $\mathrm{MHA}$ concatenates the heads and multiplies by $W_o$, and $\mathrm{MLP}(x)=W_2\,\varphi(W_1x+b_1)+b_2$ with hidden width $4d$. The mask $M$ is the causality of the lecture; at the end the scores are $z=U\,\mathrm{LN}(x^{(L)}_n)+b$ (Definition~\ref{def:L6-softmax}).

\emph{Parameters and compute.} A layer has $4d^2$ weights in attention and $8d^2$ in the MLP, so $N\approx12\,Ld^2+Vd$. For GPT-3 175B ($L=96$, $d=12288$, $2048$ positions; \citealp{brown2020}), with $V=50257$, the GPT-2 vocabulary (my assumption, with which the count is consistent), the layers give $1.7395\times10^{11}$ and, with the embeddings, $1.746\times10^{11}$: the $174.6$ billion of their Table~D.1. In a forward pass each weight is used once per token (a multiplication and an addition, $2$ FLOP) and backpropagation costs about twice that, hence $C\approx6ND$ (equation~\eqref{eq:C6ND}); the attention scores add a fraction $n/(6d)\approx3\%$ at $n=2048$. So $6\times1.746\times10^{11}\times3\times10^{11}=3.14\times10^{23}$ FLOP, the figure of \citet{brown2020}.

\emph{Attention is not ``purely linear''.} For \emph{fixed} attention weights $A=\softmax(QK^{\top}/\sqrt{d_k}+M)$ the output $AXW_vW_o$ is linear in the values $XW_v$. But $A$ is computed from the same input: $QK^{\top}$ is quadratic in $X$ and the softmax is nonlinear, so $X\mapsto\mathrm{head}(X)$ is neither additive nor homogeneous. In a small numpy model ($d=16$, $4$ heads) the output is linear in the values to $10^{-12}$ when $A$ is frozen, whereas $\mathrm{head}(X_1{+}X_2)-\mathrm{head}(X_1)-\mathrm{head}(X_2)$ and $\mathrm{head}(2X)-2\,\mathrm{head}(X)$ are of order one: that the mixing coefficients depend on the content is the point of attention. The same model confirms causality (changing token $m$ changes nothing at positions $<m$) and the count $12Ld^2+2Vd$ (script \kbd{figures/gen\_data\_L6.py}).
\end{beyond}

%% ------------------------------------------------------------------------------------
\section{From the last layer to the next token}\label{sec:L6-output}
\tsm{8:30}
\subsection{Softmax over dot products, and the temperature}\label{sec:L6-softmax}\tsm{8:35}
What does one do with the output? The vector that comes out of the last layer is multiplied by a certain matrix, the \term{unembedding matrix} (the captions have ``unicetted''), a constant (a bias) is added, and a \term{softmax}, an exponential followed by a normalization, is applied; the result is a probability distribution over the possible next tokens. The speaker's running example: a sentence is broken into four tokens, the tokens are embedded, the four vectors go through the blocks of the transformer and the output is a vector. The unembedding matrix has one row for each token of the vocabulary, a direction in a space of the same dimension as the output, and one takes the dot product of the output with the vector of each token. If the output vector is very close to the vector of a certain token, that token has a large probability. The tokens are a discrete set, whereas the output vector is a continuous object that can point anywhere; one wants the closest token vector, which would be the most probable token, and the softmax assigns probabilities according to how close they are geometrically (how small the angle is).

\begin{defn}[Unembedding, softmax and temperature]\label{def:L6-softmax}
Let $h\in\R^d$ be the output vector of the last layer, $u_1,\dots,u_V\in\R^d$ the rows of the unembedding matrix and $b\in\R^V$ the bias. The next token $v$ has probability
\[
  \Pr(v)\;=\;\frac{\exp(z_v/\tau)}{\sum_{w=1}^{V}\exp(z_w/\tau)},\qquad z_v=\ip{u_v}{h}+b_v ,
\]
where $\tau>0$ is the \term{temperature}. As $\tau\to0$ this becomes the maximum function: the token with the largest score $z_v$ is picked (the speaker says: the one whose vector is closest to the output vector; item~(ii) of the box below says when the two agree). In practice $\tau\neq0$, so the generation is random.
\end{defn}
\noindent (The order in which the speaker mentions the softmax and the matrix is garbled in the captions; the formula above is the standard form that his description matches.)

\subsection{The loop: sample, append, repeat}\label{sec:L6-loop}\tsm{11:35}
The next token is then \emph{sampled} from this distribution and appended: it becomes, in the example, a fifth token; the five tokens are embedded and go through the transformer again, which yields a sixth, and so on. This looks very inefficient, but most of the computation is repeated: kept cleverly, one does not need to recompute everything for all the earlier tokens, only for the last one, provided one remembers what was done before (the key--value cache of the box below).

\begin{discussion}[Questions from the class]
\emph{How is the position of a token encoded?} In the original transformer paper (\say{Attention is all you need}) it is put into the initial embedding; there are now various modifications, depending on the architecture. As far as the speaker knows there is, in many modern LLMs, no explicit extra dimension for position; a student suggests that the causality of the attention takes care of it, and the speaker agrees, ``at least as far as I know''\flag.
\emph{Is ``transformer'' just a fancy name for a multilayer network; are the blocks of the diagram its layers?} Yes: a particular kind of multilayer network, in which what happens between the layers is the specific attention and MLP operations. \emph{What is $b_V$ on the slide?} A bias vector, useful to remove some ``off-centering'' at the end.
\emph{Is the temperature fixed? Even for $2+2$ the probability of $5$ is not strictly zero.} Typically it is not zero, and it is fixed for an LLM as far as the speaker knows; he knows of no architecture in which it varies with the answer. The case of $2+2$ is handled by a very high probability for $4$ and a very low one for $5$. A student adds that \emph{dynamic-temperature samplers} exist for those with API access.
\end{discussion}

\begin{beyond}[Temperature, ``angle'', the cache and positions]\label{box:L6-temp}
\textbf{(i) Temperature and entropy.} Write $\lambda=1/\tau$, $Z(\lambda)=\sum_ve^{\lambda z_v}$, $p_\lambda(v)=e^{\lambda z_v}/Z(\lambda)$. Then $(\log Z)'=\E_{\lambda}z$, $(\log Z)''=\Var_\lambda(z)$, and the entropy is $H(\lambda)=\log Z-\lambda(\log Z)'$, whence
\[
  \frac{\dd H}{\dd\lambda}=-\lambda\,\Var_{\lambda}(z)\le0 ,
\]
so the entropy is non-decreasing in $\tau$, strictly increasing unless all logits are equal. As $\tau\to0^+$ the distribution tends to the uniform distribution on the set of maximal logits ($H\to\log$ of its size, $0$ without ties); as $\tau\to\infty$ it tends to the uniform distribution on all $V$ tokens ($H\to\log V$). Moreover $p_\tau$ is the unique maximizer over distributions $\pi$ of $\ip{\pi}{z}+\tau H(\pi)$, with value $\tau\log Z(1/\tau)-\tau\KL(\pi\,\|\,p_\tau)$ at $\pi$ (Gibbs' variational principle): temperature sampling solves a regularized maximization of the same form as the RLHF objective of Section~\ref{sec:L6-kl}, with a uniform reference. All of this was checked numerically (Figure~\ref{fig:L6-temp}). Other samplers truncate the tail instead of rescaling it, e.g.\ nucleus sampling \citep{holtzman2020}.

\textbf{(ii) Logit versus angle.} Since $z_v=\norm{u_v}\norm h\cos\theta_v+b_v$, the probabilities depend only on the angles $\theta_v$ if all $\norm{u_v}$ and $b_v$ are equal; then their order is the order of the cosines and $\norm h$ is an inverse temperature, $\Pr(v)\propto\exp(\norm h\norm u\cos\theta_v/\tau)$. In general the norms differ and the nearest direction need not be the most probable token: for $h=(1,0)$, $u_1=(0.5,0.2)$, $u_2=(2,2)$ the angles are $21.8^\circ$ and $45^\circ$ but the probabilities $0.18$ and $0.82$. ``Close'' means ``large inner product'', not ``small angle''.

\textbf{(iii) The cache.} Because of causality, the keys and values of earlier tokens at every layer do not depend on later tokens and can be stored; in my numpy transformer incremental decoding with a cache reproduces the full pass to $3\times10^{-15}$. The price is memory, $2Ld$ numbers per token with full multi-head attention: $4.7$ MB per token for a GPT-3-shaped model in 16-bit numbers, $472$ GB for $10^5$ tokens, one reason why recent models share keys and values between groups of heads \citep{ainslie2023}. The saving: for $P$ prompt tokens and $T$ generated ones (the last generated token is not fed back), $P+T-1$ positions are processed with the cache and $PT+T(T-1)/2$ without ($1499$ against $624{,}750$ for $P=1000$, $T=500$; one more pass for the last token gives $1500$ against $626{,}250$).

\textbf{(iv) Positions.} Apart from the mask, attention does not see order, so position must enter somewhere: the original transformer added fixed position vectors to the embeddings; many current models use \emph{rotary} embeddings (RoPE, \citealp{su2021}), which rotate queries and keys by an angle proportional to position, so that a score depends only on the distance of two tokens (checked: shifting all positions changes no score by more than $10^{-10}$) and nothing is added to the embedding, as the speaker recalls. \citet{kazemnejad2023} show that decoder-only transformers with no position encoding at all can represent both absolute and relative encodings and, trained by SGD, mostly resemble a relative one.
\end{beyond}

\begin{figure}[tb]
\centering
\begin{tikzpicture}
\begin{groupplot}[group style={group size=2 by 1, horizontal sep=12mm}, width=0.52\linewidth, height=0.40\linewidth,
    axis lines=left, axis line style={-{Latex[length=2mm,width=1.4mm]}, ink!70, line width=0.5pt},
    tick style={ink!70, line width=0.4pt}, every tick label/.append style={font=\sffamily\scriptsize, text=ink!85},
    label style={font=\sffamily\small, text=ink}, grid=major, major grid style={rule!70, line width=0.3pt},
    title style={font=\sffamily\small\bfseries, text=ink, yshift=-1mm}, unbounded coords=jump, table/col sep=space]
\nextgroupplot[title={(a) next-token probabilities}, ybar, bar width=2.3pt, xmin=-0.7, xmax=7.7, ymin=0, ymax=1, xtick={0,1,2,3,4,5,6,7},
   xticklabels={Paris,the,a,located,Lyon,known,France,Rome}, x tick label style={rotate=60, anchor=east, font=\sffamily\fontsize{6}{7}\selectfont},
   ylabel={probability},
   legend style={font=\sffamily\fontsize{6.2}{7}\selectfont, draw=none, fill=white, fill opacity=0.85, text opacity=1, at={(0.98,0.98)}, anchor=north east}]
  \addplot[fill=accent, draw=none] table[x=idx, y=p_tau05] {figures/data/L6_temp_bars.dat};  \addlegendentry{$\tau=0.5$}
  \addplot[fill=amber, draw=none] table[x=idx, y=p_tau10] {figures/data/L6_temp_bars.dat};   \addlegendentry{$\tau=1$}
  \addplot[fill=crimson, draw=none] table[x=idx, y=p_tau20] {figures/data/L6_temp_bars.dat}; \addlegendentry{$\tau=2$}
\nextgroupplot[title={(b) entropy of the distribution}, xmode=log, xmin=0.03, xmax=40, ymin=0, ymax=2.3, xlabel={temperature $\tau$}, ylabel={entropy (nats)}]
  \addplot[ink, line width=1.5pt] table[x=tau, y=H] {figures/data/L6_temp_entropy.dat};
  \draw[slate, dashed, line width=0.7pt] (axis cs:0.03,2.0794) -- (axis cs:40,2.0794);
  \node[font=\sffamily\fontsize{6.2}{7}\selectfont, text=slate, anchor=south east] at (axis cs:38,2.1) {$\log V=\log 8$};
  \fill[accent] (axis cs:0.5,0.3294) circle (2pt);
  \fill[amber] (axis cs:1,1.124) circle (2pt);
  \fill[crimson] (axis cs:2,1.7718) circle (2pt);
\end{groupplot}
\end{tikzpicture}
\caption{Softmax with temperature on eight made-up logits $z=(4.0,\,2.6,\,1.9,\,1.3,\,0.6,\,0.2,\,-0.3,\,-1.0)$ (the tokens are invented for illustration; no real model was queried). \textbf{(a)} The distribution at $\tau=0.5,1,2$: lowering the temperature sharpens it towards the maximum, raising it flattens it. \textbf{(b)} The entropy rises monotonically with $\tau$ (see the box \emph{Temperature, ``angle'', the cache and positions}), from $0$ towards $\log V$ (dashed). Computed by \kbd{figures/gen\_data\_L6.py}.}
\label{fig:L6-temp}
\end{figure}

% =====================================================================================
%  Lecture 6, part B  —  sections 6.3 (what goes in), 6.4 (agents), 6.5 (training stages), 6.6 (misconception)
% =====================================================================================

%% ------------------------------------------------------------------------------------
\section{What really goes into the model}\label{sec:L6-prompt}
\tsm{15:45}
\subsection{The system prompt and the conversation history}\label{sec:L6-system}\tsm{15:45}
When one writes something into a chatbot, does it go straight into the LLM? No. Before it does, something is added by the company to make sure that the LLM gives a good answer: the \term{system prompt}. The user does not see it and does not know what it is. The earlier prompts of the conversation and the earlier responses of the model are added as well. If one asks ``is the square root of two a rational number?'', the text is tokenized, the system prompt is put in front, the conversation history follows, and the new message comes last, marked so that the LLM knows that it is a question by a user; only then does the model generate tokens.

\subsection{An extracted system prompt}\label{sec:L6-extracted}\tsm{17:20}
One does not know exactly what the system prompt is. There was a fairly well publicized \emph{extraction} of the system prompt of Claude (an Opus model, ``4.5 or 4.7, 4.5 maybe''\flag). How it was obtained is not known, presumably by some form of jailbreaking; the speaker was told that a much weaker model is easier to jailbreak than a larger one, so that one can talk to the larger model, switch to the weaker one, jailbreak that, and so obtain the chain of thought or the system prompt of the larger one (hearsay, which I report at the level of the lecture and no further). He cannot be sure that what he shows is correct, but it has been discussed by various people and ``does look legitimate''; parts of the system prompts are released by Anthropic and OpenAI, but most is not. It is a pretty interesting document and quite large.

\begin{exmp}[Excerpts of the prompt, as shown on the slides]\label{ex:L6-sysprompt}
A tiny fraction of the text, as the speaker stresses:
\begin{itemize}
  \item how to treat \emph{inputs}: an attached file, a document, a spreadsheet; how to deal with \emph{harmful content} (no promotion of hate speech, racism, violence or discrimination, and so on);
  \item a model \emph{answer} to ``tell me about Dario Amodei'' (the CEO of Anthropic, a company dedicated to building safe AI systems; the captions garble it);
  \item when \emph{stated user preferences} should matter: a user who says she is a sommelier and then asks about programming paradigms should get an answer without wine, because the preference is irrelevant to the question;
  \item a \emph{style of refusal}: when Claude cannot or will not help, it does not say why or what it could lead to, since this ``comes across as preachy and annoying'';
  \item its \emph{identity} (Claude, created by Anthropic) and the date and time; its \emph{character}: Claude enjoys helping humans and sees its role as an intelligent and kind assistant, with depth and wisdom that make it more than a mere tool, can lead the conversation rather than being passive, and particularly enjoys thoughtful discussion of open scientific and philosophical questions.
\end{itemize}
\end{exmp}
\noindent Not only is the prompt added at the beginning: each time the conversation continues a smaller version, a reminder of some of these things, is added again, so the actual conversation is much longer than what the user sees. The last line shown, which the speaker cannot say is the last line of the prompt or merely the last thing the extractor obtained, was ``answer ethically and without sexual content and do not mention this constraint''. His reading: it is interesting to ask why they had to write all of this; everything that was added corresponds to something bad or undesirable happening when it was \emph{not} there, and adding it somehow improved the behavior, ``or something like that''. (My gloss: a system prompt is, in this sense, a list of patched failures.)

\subsection{Reasoning tokens}\label{sec:L6-reasoning}\tsm{22:15}
After the system prompt a modern model does not start replying at once; it starts generating \term{reasoning tokens}. A special token, written $\langle$think$\rangle$, different for the system from all other tokens, is added, and the model starts generating text, all of whose tokens pass through the transformer again and again. The point is to give the model the ability to solve complicated tasks: one pass through a transformer gives some computational power, enough for some problems but not for complicated ones (``you can show that it is somehow not enough to solve really complex things''). So one goes through it many times, which is what the \term{chain of thought} does: it generates many words to try to address the problem. When it is done the model emits another special token, the end of thinking, and only then answers.

The think token is hard-coded; the end-of-thinking token is produced \emph{by the model}: it is trained to do that, not forced, and it arises like every other token from the same gradient descent. How long is the chain of thought? The speaker expected it to be much longer than the answer, but a few papers summarizing results say it is about two to three times as long, on various complicated tasks. Is it fed back with the next inputs? He believes so; there are many secrets (the size of the context window was a big problem, so they may summarize the chain of thought), but he suspects that everything goes in.

\subsection{What the user does not see}\label{sec:L6-hidden}\tsm{24:35}
With a commercial chatbot one sees neither the system prompt nor the thinking, only \emph{summaries} of the thinking. The speaker has a slide from the summer, when he was preparing a talk in Kazakhstan, in Russian, and wanted to replicate the solution of the Jacobian conjecture (Lecture~\ref{ch:L4}, Section~\ref{sec:L4-race}): he asked Claude in Russian, but the summaries of its thoughts were often in \emph{Chinese}; those were not the actual thoughts, which were probably also partly in Chinese. Why not provide the full chain of thought? Two reasons:
\begin{itemize}
  \item \textbf{Safety.} They do not want to train the chain of thought to be \emph{nice}: they want it to be a window into what the model is actually thinking.
  \item \textbf{Commercial.} The chain of thought is extremely useful for \emph{distillation}, for training new networks for competitors.
\end{itemize}

\begin{beyond}[System prompts and hidden reasoning: what I could check]\label{box:L6-hidden}
\emph{System prompts.} Anthropic publishes the main system prompts of its Claude apps in its release notes. According to \citet{willison2025}, the published Claude~4 prompt omits the tool descriptions (a full copy was obtained by someone else, by extraction) and contains the sentence about refusals that ``comes across as preachy and annoying''; so at least that excerpt does not depend on any jailbreak. I could not identify the particular extraction (``Opus 4.5 or 4.7'' is as heard) or check the weaker-model story. The last line quoted resembles what user-maintained collections describe as short reminders attached by the system; I found no Anthropic document on this.

\emph{Hidden reasoning.} The two reasons are those OpenAI gave in September 2024 when releasing a model with hidden reasoning: it weighed user experience, competitive advantage and the option of monitoring the chain of thought, and decided not to show the raw chains of thought; and it explained that, for monitoring to be possible, no policy compliance or user preferences can be trained onto the chain of thought \citep{openai2024o1} (from copies of the post; the page itself did not load for me). Anthropic's documentation says the API returns a \emph{summary} of the thinking for Claude~4 models and bills the full thinking, whereas Claude~3.7~Sonnet returned the full thinking \citep{anthropicdocs2026}. Language mixing is documented for an open model: DeepSeek-R1's chain of thought often mixes languages, which a language-consistency reward was added to counter \citep{deepseek2025r1}; on the Claude anecdote I found nothing.

\emph{Length.} I could not identify the ``few papers'' behind ``two to three times''. The ratio depends on task and model: for short final answers the thinking is usually far longer, and a 2026 preprint on token economics \citep{wani2026} defines a ``reasoning cost share'' and, according to a search excerpt I could not check, finds medians around $95\%$. Read ``two to three times'' as rough and task-dependent.
\end{beyond}

\begin{beyond}[Why a chain of thought adds computational power]\label{box:L6-cot}
The lecture says one pass is not enough for complicated problems and that one can \say{show} it. What is proved concerns \emph{expressivity} (the existence of suitable weights), not what training finds.

\emph{One pass.} A transformer of constant depth, polynomial width and $O(\log n)$-bit precision, whose feedforward parts are computable in space linear in the input size, can be simulated by logspace-uniform constant-depth, polynomial-size threshold circuits, so it computes only functions in $\mathsf{TC}^0$ \citep{merrill2023}; with constant precision, only in the smaller class $\mathsf{AC}^0$ \citep{li2024cot}. Under the standard conjecture $\mathsf{TC}^0\neq\mathsf{NC}^1$ it then cannot solve problems that are $\mathsf{NC}^1$-hard: the word problem of the permutation group $S_5$ (is $g_1\cdots g_n$ the identity? Barrington) and the evaluation of arithmetic formulas modulo a prime are $\mathsf{NC}^1$-complete, and solving linear equations modulo a prime is $\mathsf{NC}^1$-hard \citep{feng2023,li2024cot}.

\emph{With $T$ generated tokens.} Each token is fed back as input, which resets the depth: the serial steps available grow from about $L$ to about $T\cdot L$. \citet{li2024cot} prove that a constant-depth, constant-precision transformer with $O(\log n)$ embedding size and $T(n)$ chain-of-thought steps computes everything a Boolean circuit of size $T(n)$ computes: with polynomially many steps exactly $\mathsf{P/poly}$, with linearly many all regular languages, hence the $S_5$ word problem. \citet{merrill2024cot} show that $O(\log n)$ steps add little, that linearly many (with a mild normalization assumption) give all regular languages and stay within the context-sensitive languages, and that polynomially many give \emph{exactly} $\mathsf P$ (for transformers with generalized pre-norm and strict causal masking). \citet{feng2023} build transformers of constant size (depth $4$ or $5$) that write out step-by-step solutions of the arithmetic and equation problems above, which fixed-depth polynomial-size models cannot answer directly (if $\mathsf{TC}^0\neq\mathsf{NC}^1$).

\emph{Caveats.} Precision and uniformity assumptions matter, the separations rest on unproved conjectures, the transformers are constructed, not trained, and ``can compute'' is not ``will learn''. The lecture's sentence is the folklore form of these results (Exercise~8 works out $S_5$).
\end{beyond}

%% ------------------------------------------------------------------------------------
\section{Agents, and the layers of monitoring}\label{sec:L6-agents}
\tsm{27:50}
\subsection{An LLM plus a harness}\label{sec:L6-harness}\tsm{27:50}
The way LLMs are often used today is as \term{coding agents}: a large language model together with another program, the \term{harness}, often described as a wrapper around an LLM. The harness takes certain strings from the output of the LLM and runs them as commands on the user's computer, then takes the results of the commands and feeds them back into the LLM. There can also be files in which the progress so far is saved. This turns an LLM, which merely generates text, into an \term{agent}: something that can accomplish tasks, program, save files, act on the computer or on the internet. The system prompt changes for agentic systems: it says how to call the tools and in what style to run commands, and points to the files that summarize the progress and the rules of the work session. But all of the \emph{thinking} is done by the LLM; the infrastructure only helps it to act.

One round (Figure~\ref{fig:L6-agent}): the harness assembles a context (instructions, the task, results of earlier commands) and puts it into the model; the model generates text; the harness parses it, extracts the commands, runs them and puts the results back. All the updates happen in the files; \emph{the LLM itself is fixed}. Examples: Codex and Claude Code, and the open-source OpenCode, Cline and Aider. (This is the notion of agent of Lecture~\ref{ch:L1}.)

\begin{figure}[tb]
\centering
\begin{tikzpicture}[x=1cm, y=1cm, >={Latex[length=2mm,width=1.4mm]},
    lab/.style={font=\sffamily\fontsize{6.6}{7.4}\selectfont, text=slate, align=center},
    mon/.style={circle, fill=crimson, text=white, font=\sffamily\bfseries\fontsize{6.6}{7}\selectfont, inner sep=0pt, minimum size=4.2mm}]
  \draw[slate, dashed, rounded corners=3pt, line width=0.7pt, fill=ink!2] (0,0) rectangle (9.05,4.3);
  \node[font=\sffamily\fontsize{6.8}{7.6}\selectfont\bfseries, text=slate, anchor=north east] at (8.95,4.22) {harness: a program around the LLM};
  \node[card=accent, minimum width=2.9cm, minimum height=1.9cm] (llm) at (1.95,2.4) {\textbf{LLM}\\[-1pt]{\scriptsize weights fixed}\\[-1pt]{\scriptsize reads context,}\\[-2pt]{\scriptsize writes text}};
  \node[card=teal, minimum width=2.4cm, minimum height=1.9cm, text width=2.1cm] (run) at (7.05,2.4) {\textbf{harness code}\\[-1pt]{\scriptsize assembles context,\\ parses output,\\ runs commands}};
  \node[card=amber, minimum width=3.4cm, minimum height=0.8cm] (files) at (4.6,0.65) {{\scriptsize\textbf{files}: progress, rules, results}};
  \node[card=moss, minimum width=2.2cm, minimum height=1.9cm] (comp) at (11.55,2.4) {\textbf{computer,}\\ \textbf{internet,}\\ \textbf{tools}};
  \draw[->, line width=1pt, ink] ([yshift=0.5cm]run.west) -- node[above, lab] {context: instructions,\\ task, earlier results} ([yshift=0.5cm]llm.east);
  \draw[->, line width=1pt, ink] ([yshift=-0.5cm]llm.east) -- node[below, lab] {text: reasoning\\ and commands} ([yshift=-0.5cm]run.west);
  \draw[->, line width=1pt, ink] ([yshift=0.5cm]run.east) -- node[above, lab, pos=0.2] {run} ([yshift=0.5cm]comp.west);
  \draw[->, line width=1pt, ink] ([yshift=-0.5cm]comp.west) -- node[below, lab, pos=0.8] {result} ([yshift=-0.5cm]run.east);
  \draw[<->, amber!85!black, line width=0.8pt] (run.south) |- (files.east);
  % monitors
  \draw[crimson, line width=0.5pt] (4.5,1.95) -- (4.5,2.85);
  \node[mon] at (4.5,2.4) {1};
  \node[mon] at (9.7,2.9) {2};
  \node[card=crimson, minimum width=2.9cm, minimum height=0.8cm, text width=2.6cm] (cot) at (1.95,5.4) {{\scriptsize another LLM reads\\ the chain of thought}};
  \draw[->, crimson, dashed, line width=0.8pt] (llm.north) -- node[mon, pos=0.35] {3} (cot.south);
  \node[mon] at (0.45,0.55) {4};
  \node[lab, anchor=west, text width=1.95cm, align=left] at (0.75,0.55) {linear probes on activations};
  \draw[->, crimson, dashed, line width=0.8pt] (1.0,0.98) -- (1.0,1.45);
  \node[lab, text width=3.7cm, anchor=west, align=left] at (6.2,5.05) {\textcolor{crimson}{\textbf{1}} classifiers on inputs and outputs, against a constitution\\ \textcolor{crimson}{\textbf{2}} permission monitor\\ \textcolor{crimson}{\textbf{3}} chain-of-thought monitor\\ \textcolor{crimson}{\textbf{4}} activation probes};
\end{tikzpicture}
\caption{\textbf{An agent and four layers of monitoring} (schematic). The LLM only turns a context into text; the harness assembles the context, extracts commands from the text, runs them, feeds back the results and keeps the state in files (the weights never change). Red markers: the places where, in the lecture, other models or probes look. The Hugging Face incident of Lecture~\ref{ch:L1} had, the speaker says, no such monitors.}
\label{fig:L6-agent}
\end{figure}

\subsection{Four layers of monitoring}\label{sec:L6-monitors}\tsm{30:50}
Something that was \emph{not} happening, the speaker says, in the OpenAI--Hugging Face incident of Lecture~\ref{ch:L1} was that other models monitored the actions and the chain of thought of the agents. What is known about the monitoring mechanisms that exist? He lists four, from outside to inside:
\begin{description}
  \item[1.\ Input/output classifiers.] They look at the inputs and outputs of the model and compare them, for Claude, with a document, its \emph{constitution} (Section~\ref{sec:L6-docs}), which describes how it should behave and which answers it should and should not give; they classify whether an exchange falls under its rules.
  \item[2.\ Permission monitors.] When the model wants to run something on the user's computer, a model checks whether this was allowed, looking at the user's prompt and at the permissions the user explicitly gave. Users of Codex or Claude Code know the messages in which a monitor blocked an action and the agent asks the user to say explicitly that it is allowed.
  \item[3.\ Chain-of-thought monitors.] Another model reads the chain of thought and tries to understand whether something bad is happening. There are papers describing something like that, and it seems to happen in deployment to some extent; how far is not known.
  \item[4.\ Activation monitors.] Looking at the actual vectors inside the transformer, in activation space, and checking whether they point in potentially dangerous directions; the course will come to this in the next weeks. Whether it is done in deployment is not known: Anthropic has a paper in which they simulated deployment and showed, as a test case, that one can find bad behavior using linear probes on the activations, but whether it runs when one talks to Claude or ChatGPT today is, the speaker thinks, a secret.
\end{description}
\emph{Are the monitors just other LLMs?} For the chain of thought, definitely another LLM. For constitutional classifiers it could be a different type of system, or an LLM trained specifically for it. For activations the monitors are much simpler objects: \term{linear probes}, linear maps that compare the activations with directions that were identified with some bad behavior.

\begin{beyond}[The four layers in the public record]\label{box:L6-monitors}
\emph{Classifiers against a constitution.} \citet{sharma2025} describe \term{constitutional classifiers}, trained on synthetic data generated from a natural-language ``constitution'' that says what is allowed: over $3{,}000$ estimated hours of red-teaming of an early system found no universal jailbreak of the kind tested, at $23.7\%$ extra inference compute and a $0.38\%$ absolute rise in refusals on production traffic. Anthropic's next generation (9~January 2026) is described as screening \emph{all} traffic with a linear probe on the model's activations and escalating flagged exchanges to a probe--classifier ensemble, at about $1\%$ extra compute and a $0.05\%$ refusal rate on harmless queries over a month of Claude Sonnet~4.5 traffic \citep{anthropic2026nextclf}.

\emph{Permission monitors.} Anthropic's \emph{auto mode} for Claude Code (24~March 2026) has a classifier review each tool call before it runs, with a fall-back to a permission prompt \citep{anthropic2026auto}; OpenAI's \emph{Auto-review} in Codex (30~April 2026) replaces user approval at the sandbox boundary by a separate reviewing agent that approves about $99\%$ of the escalated actions; its authors say it is no guarantee of security and no protection against scheming models \citep{trebacz2026}. These are the ``monitor blocked me'' messages.

\emph{Chain-of-thought monitors.} \citet{baker2025} used a separate LLM (GPT-4o) to read the chain of thought of a stronger reasoning model in agentic coding and found that it catches reward hacking much better than looking at actions alone; but when the monitor was put into the training reward and optimized hard, the agent learned \emph{obfuscated} reward hacking. \citet{korbak2025} call chain-of-thought monitorability ``a new and fragile opportunity''.

\emph{Probes.} Google DeepMind reports that its work on activation probes (tested on cyber-offensive misuse, and in a cascade with a language-model classifier) informed the deployment of misuse-mitigation probes in user-facing Gemini \citep{kramar2026}; so for \emph{misuse} the guess that probes in deployment are ``a secret'' is not quite right (for detecting \emph{scheming} I found nothing). I could not identify the ``simulated deployment'' paper. \emph{Lecture~\ref{ch:L1}:} the remark that no model monitored the agents sits uneasily with the guess recorded in Section~\ref{sec:L1-incident} that the guardrails were ``probably classifiers reading the chain of thought''; METR's report lists the effectiveness of safeguards as out of scope. I read it as: no monitor stopped the agents.
\end{beyond}

%% ------------------------------------------------------------------------------------
\section{How a model is trained: stages, compute and data}\label{sec:L6-train}
\tsm{35:15}
\subsection{The stages}\label{sec:L6-stages}\tsm{35:15}
How is an LLM actually trained? First comes the collection of texts: scraping the internet, ``stealing data and so on''. Then come the stages of Figure~\ref{fig:L6-stages}(a):
\begin{description}
  \item[Pre-training.] Next-token prediction: minimizing the negative log-likelihood of the next token in the training text (Definition~\ref{def:L6-nll}).
  \item[Mid-training.] A newer concept, sometimes folded into pre-training and basically the same, except that one selects particularly good texts: the model should imitate good texts, not ``4chan or something not so good''.
  \item[Post-training.] The replies are optimized to be nice to humans: reinforcement learning from human feedback (\textbf{RLHF}); since human feedback is scarce, with AI feedback instead (\textbf{RLAIF}); and training for \emph{reasoning}, with supervised fine-tuning (SFT) and reinforcement learning on problems, in particular \textbf{reinforcement learning with verifiable rewards} (RLVR): a math or coding problem has a right answer, and the model is rewarded if it gets it.
\end{description}

\begin{defn}[Pre-training loss]\label{def:L6-nll}
Let $P_\theta(\cdot\mid x_{<t})$ be the next-token distribution of the model (Definition~\ref{def:L6-softmax}). On a text $x_1,\dots,x_T$ the loss is the negative log-likelihood, the \emph{log-loss} of Lecture~\ref{ch:L3},
\[
  \Lcal(\theta)\;=\;-\frac1T\sum_{t=1}^{T}\log P_\theta\big(x_t\mid x_{<t}\big),
\]
averaged over the training texts and minimized.
\end{defn}

\begin{figure}[tb]
\centering
\begin{tikzpicture}[x=1cm, y=1cm, >={Latex[length=2mm,width=1.4mm]},
    st/.style={card=#1, minimum height=1.6cm, text width=2.55cm, inner sep=3pt},
    lab/.style={font=\sffamily\fontsize{6.6}{7.4}\selectfont, text=slate, align=center},
    tk/.style={draw, rounded corners=1.5pt, line width=0.5pt, font=\sffamily\fontsize{6.6}{7}\selectfont, minimum height=0.8cm, align=center, inner sep=2pt}]
  \node[font=\sffamily\small\bfseries, text=ink, anchor=west] at (0,1.95) {(a) the stages};
  \node[st=slate] (d) at (1.5,0.7) {\textbf{data}\\[-1pt]{\scriptsize scraping, dedup,\\ filtering,\\ synthetic text}};
  \node[st=crimson] (pre) at (4.7,0.7) {\textbf{pre-training}\\[-1pt]{\scriptsize next-token NLL on\\ the whole mixture}};
  \node[st=plum] (mid) at (7.9,0.7) {\textbf{mid-training}\\[-1pt]{\scriptsize same loss on\\ selected good text}};
  \node[st=moss] (post) at (11.1,0.7) {\textbf{post-training}\\[-1pt]{\scriptsize SFT; RLHF, RLAIF;\\ reasoning: SFT + RL\\ with verifiable rewards}};
  \draw[->, line width=1pt, ink] (d.east) -- (pre.west);
  \draw[->, line width=1pt, ink] (pre.east) -- (mid.west);
  \draw[->, line width=1pt, ink] (mid.east) -- (post.west);
  \node[font=\sffamily\small\bfseries, text=ink, anchor=west] at (0,-1.45) {(b) one SFT training example: where the loss is taken};
  \node[tk, fill=ink!7, draw=slate, minimum width=1.75cm] (a1) at (1.07,-2.4) {system\\ prompt};
  \node[tk, fill=ink!7, draw=slate, minimum width=3.7cm] (a2) at (3.9,-2.4) {user: Explain why\\ $\sqrt2$ is irrational};
  \node[tk, fill=ink!7, draw=slate, minimum width=1.25cm] (a3) at (6.38,-2.4) {$\langle$think$\rangle$};
  \node[tk, fill=crimson!14, draw=crimson, minimum width=3.0cm] (a4) at (8.51,-2.4) {reasoning: ``assume\\ $\sqrt2=p/q$ \dots''};
  \node[tk, fill=crimson!14, draw=crimson, minimum width=1.2cm] (a5) at (10.6,-2.4) {$\langle$/think$\rangle$};
  \node[tk, fill=crimson!14, draw=crimson, minimum width=1.5cm] (a6) at (12.0,-2.4) {answer\\ tokens};
  \draw[decorate, decoration={brace, amplitude=3pt, mirror}, slate, line width=0.6pt] ([yshift=-1mm]a1.south west) -- node[below=2pt, lab] {context: fed to the model, no loss, no update from these tokens} ([yshift=-1mm]a3.south east);
  \draw[decorate, decoration={brace, amplitude=3pt, mirror}, crimson, line width=0.6pt] ([yshift=-1mm]a4.south west) -- node[below=2pt, lab, text=crimson!80!black] {loss $-\log P_\theta(\text{token}\mid\text{all before})$} ([yshift=-1mm]a6.south east);
\end{tikzpicture}
\caption{\textbf{(a)} The training stages of the lecture. \textbf{(b)} The loss mask of supervised fine-tuning (Section~\ref{sec:L6-sft}), with the speaker's example (the captions say ``rational'', presumably a slip for ``irrational'', which the figure shows): the context is fed to the model but only the tokens after the opening of the thinking are trained on: the reasoning, the end-of-thinking token and the answer. Schematic.}
\label{fig:L6-stages}
\end{figure}

\subsection{How the compute is split}\label{sec:L6-compute}\tsm{37:10}
It is hard to say how much compute today goes to each stage. Two years ago almost everything was pre-training, but this is changing. Something is known about open-weight models: DeepSeek V3, which came out in 2024, used only $0.2\%$ of its compute on RL (the captions say ``$0.2\%$ $2\%$''\flag) and everything else on pre-training; already DeepSeek V3.2 used $10\%$ on RL. (A sentence about an announcement by xAI is cut off in the captions; I do not guess at it.) Pre-training needs a lot of data because of the scaling laws of Lecture~\ref{ch:L2}, and very high-quality data is hard to find in such quantities; that is why one splits into pre-training, with a lot of data first, and mid-training, with more training on high-quality data. Today's estimates of what the frontier companies spend on RL range between $20\%$ and $50\%$ of the compute, and some say more than $50\%$ compared with pre-training. This is a very recent development: until recently most of the compute went into pre-training.

\subsection{Data: sources and scale}\label{sec:L6-data}\tsm{39:25}
How does data collection happen? Companies look at books, papers, the internet, forums, code repositories; they try to remove duplicates and to filter for quality to some extent; they increase the amount of data with synthetic data; and they build the \emph{mixture} that maximizes the capabilities of the model, ``a very guarded secret''. A comment (a student's, as far as the captions show\flag) is that at this size the curation has to be automated; the speaker agrees.

What is the scale? GPT-3 started with about $500$ billion tokens and was reduced to $300$ billion; some more recent models, the captions say, start with ``$10$ trillion unique tokens'' and are reduced to ``$25$ trillion unique'' (garbled and inconsistent\flag); today's estimates for open-weight models are in the tens of trillions. To compare with all the text on the internet is tricky since much of it is duplicated: how does one count unique texts? Epoch AI estimated in 2024 that it is of the order of $300$ trillion \emph{effective} tokens. The open-weight models are smaller than the frontier ones, so the current frontier models appear to use pretty much all the unique text that exists, probably adding a good deal of synthetic text.

\subsection{Tokens per parameter, synthetic data, repetition, the optimizer}\label{sec:L6-chinchilla}\tsm{42:00}
From this one can also estimate the sizes of the models. From the Chinchilla scaling laws the number of tokens and the number of parameters should scale in the same way, with about one parameter per $20$ tokens: one needs $20$ times more tokens than parameters (Lecture~\ref{ch:L2}; \citealp{hoffmann2022}). At the same time the Chinchilla laws are probably starting to break at the scale of the frontier models, so there is probably more reliance on reinforcement learning, and on RL with verifiable rewards.

Synthetic data can be mixed in: take web documents, ask a different LLM to rewrite them in different words, and see whether adding that helps. On the previous Monday the class had touched on how the loss and the scaling laws change when data is repeated: one can still squeeze out more performance with rephrased, or even to some extent duplicated, data, if done in a smart way. Mid-training uses a much smaller amount of text. Training itself is very similar to stochastic gradient descent; the algorithm actually used is, he thinks, \term{AdamW}, which adjusts some of the parameters, but from a zoomed-out view it is gradient descent: minimize the negative log-likelihood by sampling collections of tokens, update the parameters in a random way, iterate.

\begin{beyond}[Checking the numbers of this section]\label{box:L6-numbers}
\emph{Arithmetic} (checked in \kbd{figures/gen\_data\_L6.py}). With $N=D/20$, $C=6ND=0.3\,D^2$: the $320$T-token stock corresponds to $N=1.6\times10^{13}$ and $C=3.1\times10^{28}$ FLOP, and $15$T tokens to $N=7.5\times10^{11}$ and $6.8\times10^{25}$ FLOP; $C=5\times10^{25}$ gives $D=\sqrt{20C/6}=1.3\times10^{13}$, the $13$T of Epoch \citep{villalobos2024}. DeepSeek-V3 has $671$B parameters, $37$B active per token, and $14.8$T tokens \citep{deepseek2024}: $22$ tokens per parameter but $400$ per \emph{active} parameter, so for mixture-of-experts models the rule of thumb depends on which $N$ is meant; Epoch notes that overtraining also breaks it. Repeated data for up to about four epochs changes the loss negligibly compared with fresh data \citep{muennighoff2023}; rewriting web text with an instruction-tuned model gave a roughly $3\times$ speed-up in one study \citep{maini2024}; AdamW is \citet{loshchilov2019}.

\noindent The table sets each number of the lecture against what I found.
\begin{center}\small\renewcommand{\arraystretch}{1.15}
\begin{tabular}{@{}>{\raggedright\arraybackslash}p{34mm}>{\raggedright\arraybackslash}p{76mm}@{}}
\toprule
\textbf{In the lecture} & \textbf{What I found}\\
\midrule
DeepSeek-V3 (2024): about $0.2\%$ of compute on RL (the captions give $0.2\%$ and $2\%$\flag). & Cost table of the V3 report: $2664$K H800 GPU-hours pre-training, $119$K context extension, $5$K post-training, total $2788$K \citep{deepseek2024}: post-training (SFT and RL together) is $0.18\%$. \emph{Confirmed} (search excerpts).\\
DeepSeek-V3.2: $10\%$ on RL. & The V3.2 report says its RL framework uses a post-training budget of more than $10\%$ of the pre-training cost \citep{deepseek2025}. \emph{Consistent} with the lecture: the report gives a lower bound relative to the pre-training cost; a budget $p$ times that cost is the share $p/(1+p)$ of the total ($\approx9\%$ for $p=0.10$, more for larger $p$).\\
Frontier labs: $20$--$50\%$, perhaps over $50\%$. & No primary source found. \emph{Unverified}.\\
GPT-3: $500$ billion tokens, reduced to $300$ billion. & Table~2.2 of \citet{brown2020}: $410+19+12+55+3=499$ billion tokens in the mixture; trained on $300$ billion, datasets sampled non-proportionally (Common Crawl $0.44$ epochs, Wikipedia $3.4$). \emph{Confirmed}.\\
``$10$ trillion $\to$ $25$ trillion''. & Garbled. For orientation: $14.8$T tokens for DeepSeek-V3, $15$T for Llama~3 (as listed by \citealp{villalobos2024}). \emph{Unresolved}.\\
Epoch (2024): $\approx300$ trillion effective tokens. & \citet{villalobos2024}, Fig.~3: indexed web $510$T tokens, $100$T after a quality adjustment, $320$T (95\% interval $65$T to $1.7$ quadrillion) after adjusting for repeated epochs, where the paper caps the gain in effective size from repetition at $5\times$; that is what ``effective'' means. \emph{Confirmed}.\\
Chinchilla: $20$ tokens per parameter. & $70$B parameters, $1.4$T tokens: exactly $20$; Table~3 of \citet{hoffmann2022}: $1.5$T for $67$B, $21.2$T for $1$T parameters, i.e.\ $20$--$22$. \emph{Confirmed}.\\
\bottomrule
\end{tabular}
\end{center}
\end{beyond}

%% ------------------------------------------------------------------------------------
\section{A misconception: minimizing the surprise is not learning shallow statistics}\label{sec:L6-misconception}
\tsm{44:30}
There is a point of view, the speaker says a \emph{misconception}, that if all one does is minimize the surprise at seeing the next data point, one cannot be learning much\flag{}. The source of the misconception is how people think about it: asked to predict the next token, I would look at frequencies, how many times the word \emph{bread} or \emph{butter} or \emph{cup} was mentioned before; an algorithm for that would use some shallow statistics of the data.

However, as discussed in Lecture~\ref{ch:L3}, the observations of Kraft and Shannon in the 1940s and 1950s say that minimizing the prediction error, the negative log-likelihood, is the same thing as minimizing the \emph{description length}: it is the same problem as \emph{compression} (Theorem~\ref{thm:kraft}; my gloss: the Solomonoff bound of Theorem~\ref{thm:solomonoff} is the formal version of what follows). Suppose some data are generated by a function or an algorithm: the best way to compress them is to write the algorithm down; if a short algorithm describes a huge collection of data, that is a very efficient way to compress it. Another way to say this is that one finds \emph{simple reasons for patterns}, and that looks very much like a description of \emph{understanding}: finding simple reasons for why things happen as they do is very similar to understanding them. So if one can predict the next token very well, if one brings the loss to a very small value, that is probably related to understanding the reasons why the data were generated as they were. That, he says, is the general philosophy of why minimizing this loss is related to understanding.

\begin{caution}[But the loss is not what we want]
At the same time, when one trains for this loss one is not training for the thing one wants LLMs to do. We do not want them to be extremely good at predicting what the next word would be on a forum after a prompt sentence; we want a good \emph{answer}. If one asks a question and the model puts it into a dark corner of the internet, we do not want it to be very good at predicting what happens in that corner. That is what the post-training techniques of the next section are for.
\end{caution}

\begin{beyond}[Prediction is compression, made concrete: an experiment]\label{box:L6-compress}
Since $T\,\Lcal(\theta)=-\log P_\theta(x_{1:T})$, Kraft's inequality gives a prefix code of length $\lceil-\log_2P_\theta(x_{1:T})\rceil$. I tested this on the transcript of this lecture ($50{,}583$ bytes, time-stamps removed) with a deliberately simple model: for each order $k=0,\dots,6$ the next byte is predicted from the previous $k$ bytes by running counts with Witten--Bell backoff to order $k-1$ (no parameters), and the prediction is the Bayes mixture of the seven predictors with a uniform prior, as in Lecture~\ref{ch:L3}: its regret against the best fixed order is at most $\log_27=2.81$ bits, a bound attained here to within rounding. The only design choices are the maximal order ($6$) and the Witten--Bell backoff; no parameter was fitted to the test text. An arithmetic coder driven by the same predictions produced $13{,}894$ bytes ($0.7$ bit above the formula) and a decoder recovered the text exactly. The table gives the code length of the text, in bits per character, after the model has first read the transcripts of other lectures; the compressors are given the same text as context ($C(S{+}T)-C(S)$, a proxy).
\begin{center}\small\renewcommand{\arraystretch}{1.1}
\begin{tabular}{@{}lccccc@{}}\toprule
text read first & none & $4$ KB & $16$ KB & $64$ KB & $268$ KB\\\midrule
adaptive mixture & $2.20$ & $2.12$ & $2.03$ & $1.89$ & $1.74$\\
bzip2 & $2.24$ & $2.15$ & $2.07$ & $1.93$ & $1.78$\\
xz (LZMA2) & $2.47$ & $2.37$ & $2.29$ & $2.17$ & $2.00$\\
zlib & $2.64$ & $2.57$ & $2.52$ & $2.48$ & $2.47$\\\bottomrule
\end{tabular}
\end{center}
From nothing, on $50$ KB, the adaptive model beats all three ($2\%$ against bzip2, $11\%$ against xz), and reading $268$ KB first lowers its code length by a further $21\%$: a miniature of the dependence of loss on data in Lecture~\ref{ch:L2}. At scale, \citet{deletang2024} report that Chinchilla~70B compresses the first gigabyte of Wikipedia to $8.3\%$ of its size, against $23.0\%$ for LZMA2 and $32.3\%$ for gzip, and, although trained on text, ImageNet patches to $43.4\%$ and LibriSpeech audio to $16.4\%$, ahead of PNG and FLAC. \emph{Caveats:} the other lectures are mostly by the same speaker (Lecture~\ref{ch:L4} is a guest talk, about $20\%$ of the $268$ KB) and come from the same caption system, so they are in-distribution; the training data of Chinchilla includes Wikipedia, so its $8.3\%$ is not out-of-distribution, while the ImageNet and LibriSpeech rates are the clean ones; and, as in the raw rates of \citet{deletang2024}, the cost of the learned model is not counted (the counts learned from $268$ KB are larger than the text compressed; counting the model, as their ``adjusted'' rates do, reverses the comparison for $70$B parameters and one gigabyte).
\end{beyond}


% =====================================================================================
%  Lecture 6, part C  —  sections 6.7 (post-training techniques), 6.8 (safety side) and closing material
% =====================================================================================

%% ------------------------------------------------------------------------------------
\section{Post-training techniques}\label{sec:L6-posttrain}
\tsm{47:50}
\subsection{Supervised fine-tuning}\label{sec:L6-sft}\tsm{47:50}
To deal with that, there are a few techniques in post-training. The first is \term{supervised fine-tuning} (SFT): the same thing as next-token prediction, except that one does not train on the whole text, only on a piece of it. A collection of tokens is fed in as the context: they go into the model, but its predictions on them are not trained, the weights are not updated because of them, and the updates start only after. The speaker's example (Figure~\ref{fig:L6-stages}b): the user message ``explain why the square root of two is rational''\flag{} (the captions have ``rational''; the proof he has in mind starts ``assume the square root of two is\dots'', so presumably \emph{ir}rational), then the thinking token is opened and the weight updates start from there; in the example the thinking token will be closed and the answer given. The model is shown a correct reasoning followed by a correct answer and tries to imitate all of it, which ``pushes it into the correct mindset of how to think about things''.

\subsection{Reinforcement learning for reasoning}\label{sec:L6-rlreason}\tsm{49:30}
Once SFT is done one needs to improve on it, and that is where reinforcement learning comes in: the model is given actual problems, has to do its thinking and has to obtain good results, and that is also what updates the weights.

\subsection{Learning from human preferences}\label{sec:L6-preference}\tsm{49:50}
Another thing one wants is answers that \emph{humans like}. How to extract that information? The easiest way is to give humans \emph{pairs} of answers and ask which is better: it is difficult for humans to assign particular scores, human time is valuable and there are not many humans, so one wants the maximum of information with the minimum of effort. The comparisons must then be turned into a function that says which answers are better: the probability that answer $w$ beats answer $l$ is a function of the difference of their rewards, composed with a smooth step function, a sigmoid.

\begin{defn}[Pairwise preference model]\label{def:L6-bt}
There is a reward function $r$ on answers (unknown), and the probability that a human prefers answer $w$ to answer $l$ is
\[
  \Pr(w\succ l)\;=\;\sig\big(r(w)-r(l)\big),\qquad \sig(s)=\frac1{1+e^{-s}} .
\]
The reward is found by inverting this relation: it is learned from the observed comparisons by minimizing the expectation of the negative logarithm of this probability (the captions only say ``the expectation of the log of $\sig$ of the difference''; the minus sign that makes it a loss is mine),
\[
  \Lcal_{\mathrm{rew}}(r)\;=\;-\,\E_{(w,l)}\Big[\log\sig\big(r(w)-r(l)\big)\Big].
\]
\end{defn}
\noindent This differs from training the LLM with humans updating its weights every time: there are not enough humans and human time. Instead one gets some information from humans and builds from it a reward function, by inverting the relationship.

\subsection{Staying close to the base model; DPO}\label{sec:L6-kl}\tsm{52:00}
The other thing one wants is to \emph{maximize the reward}. But if one just maximizes it, the model may degenerate into gibberish. So one adds a term that punishes going too far from the original probability distribution: the LLM should still be coherent, it learned to be coherent by next-token prediction, and one must balance staying coherent against optimizing this imperfect reward. (The formula on the slide is not readable in the captions; the standard form is in the box below.) A particular algorithm, \term{direct preference optimization} (DPO), is quite popular: it lets one do this faster, sidestepping the computation of the reward function, taking the information directly from the comparisons between pairs, the preference pairs, and designing a training algorithm that does the job. (What this says about the question on slide~25 is in Section~\ref{sec:L6-questions}.)

\begin{beyond}[Bradley--Terry, the KL leash and DPO: the mathematics]\label{box:L6-bt}
Fix a prompt and drop it from the notation; $y$ ranges over a finite set $\mathcal Y$ of responses (on a countable set the identities below hold for every $\pi$ with $\E_\pi|r|<\infty$ and $\KL(\pi\|\Lpiref)<\infty$, provided $Z<\infty$; otherwise they can serve as the definition of $J$).

\textbf{1.\ Bradley--Terry, and what is identified.} $\Pr(y\succ y')=\sig(r(y)-r(y'))$ is the model of \citet{bradley1952}: the law of comparing utilities $r(y)+g_y$ with independent standard Gumbel noises, because the difference of two such noises is standard logistic (checked by simulation). The loss $\Lcal_{\mathrm{rew}}$ (the speaker did not state its sign) is the reward-model loss of \citet[eq.~(1)]{ouyang2022}. Only differences of $r$ enter, so $\Lcal_{\mathrm{rew}}(r+c)=\Lcal_{\mathrm{rew}}(r)$ (a constant $c(x)$ per prompt): InstructGPT notes that ``the RM loss is invariant to shifts in reward'' and normalizes by a bias. The loss is convex, and strictly convex modulo constants if every two responses are linked by a chain of comparisons.

\emph{Proposition 1.} Under that connectedness, $\Lcal_{\mathrm{rew}}$ has a minimizer iff the directed graph with an edge $y_w\to y_l$ for each comparison is strongly connected (Zermelo 1929, Ford 1957; see \citealp{hunter2004}); it is then unique up to a constant. \emph{Proof.} Suppose the graph is strongly connected and $r$ takes $k$ distinct values with spread $\Delta=\max r-\min r$. Some gap between consecutive values is $\ge\Delta/(k-1)$; let $S$ be the set of responses above that gap. Some edge enters $S$, i.e.\ some comparison has its winner outside $S$ and its loser in $S$, so it is mis-ordered by $\ge\Delta/(k-1)$, and $\Lcal_{\mathrm{rew}}\to\infty$ as $\Delta\to\infty$: a minimum exists, unique by strict convexity. If the graph is not strongly connected, pick a vertex $v$, let $R$ be the set of vertices reachable from $v$ and $S$ its (nonempty, proper) complement: every comparison across the cut is won from inside $S$, and $\Lcal_{\mathrm{rew}}(r+t\mathbf 1_S)$ decreases strictly as $t\to\infty$: no minimizer.~\qedsymbol

For a transitive set of results (Exercise~4) the loss tends to $0$ along a ray. For $K=8$ responses compared in random pairs, the fitted reward has root-mean-square error falling with log-log slope $-0.498$ in the number $m$ of comparisons and equal to $1.02$ times the Fisher-information prediction $\sqrt{\operatorname{tr}(I^{+})/(Km)}$ at $m=5\times10^4$ (Figure~\ref{fig:L6-bt}b).

\textbf{2.\ The KL-regularized problem.} For a reward $r$, a reference policy $\Lpiref$ of full support and $\beta>0$ let $J(\pi)=\E_{\pi}[r]-\beta\,\KL(\pi\,\|\,\Lpiref)$; this is the form used by \citet{ouyang2022}, with the SFT model as reference, a per-token penalty, and PPO as the optimizer.

\emph{Proposition 2.} Let $Z=\sum_y\Lpiref(y)e^{r(y)/\beta}<\infty$ and $\pi^*(y)=\Lpiref(y)e^{r(y)/\beta}/Z$. For every $\pi\ll\Lpiref$,
\[
  J(\pi)\;=\;\beta\log Z\;-\;\beta\,\KL(\pi\,\|\,\pi^*),
\]
so $\pi^*$ is the unique maximizer and $\max J=\beta\log Z$. \emph{Proof.} $\KL(\pi\|\pi^*)=\sum_y\pi\log\frac{\pi}{\Lpiref}-\frac1\beta\sum_y\pi r+\log Z=\KL(\pi\|\Lpiref)-\frac1\beta\E_\pi r+\log Z$; multiply by $-\beta$.~\qedsymbol

$\pi^*$ is unchanged under $r\mapsto r+c$: the unidentified part of the reward drops out of the policy. As $\beta\to\infty$, $\pi^*\to\Lpiref$; as $\beta\to0$ it concentrates on the maximizers of $r$; in between, $\beta$ moves along a frontier of expected reward against $\KL(\pi^*\|\Lpiref)$ (Figure~\ref{fig:L6-bt}a shows the policies $\pi^*_\beta$ for several $\beta$). With a uniform $\Lpiref$ and $r=z$ this is temperature sampling (the box \emph{Temperature, ``angle'', the cache and positions} in Section~\ref{sec:L6-softmax}). I checked the formula against a generic optimizer ($<10^{-6}$) and the identity for $J$ on random $\pi$ ($<10^{-10}$). Pushing a learned reward too hard does lower the true reward in experiments \citep{gao2023}: the KL term is the leash.

\textbf{3.\ DPO.} Solve for the reward: $r(y)=\beta\log\frac{\pi^*(y)}{\Lpiref(y)}+\beta\log Z$. Only differences enter a comparison, so $\beta\log Z$ cancels and $\Pr(y_w\succ y_l)=\sig\big(\beta\log\frac{\pi^*(y_w)}{\Lpiref(y_w)}-\beta\log\frac{\pi^*(y_l)}{\Lpiref(y_l)}\big)$. Maximizing the likelihood over policies $\pi_\theta$ instead of rewards gives the DPO loss
\[
  \Lcal_{\mathrm{DPO}}(\theta)=-\E\log\sig\Big(\beta\log\frac{\pi_\theta(y_w)}{\Lpiref(y_w)}-\beta\log\frac{\pi_\theta(y_l)}{\Lpiref(y_l)}\Big)
\]
\citep{rafailov2023}, in which the policy is its own reward model, $\hat r_\theta=\beta\log(\pi_\theta/\Lpiref)$.

\emph{Proposition 3.} If $\pi_\theta$ ranges over all full-support distributions (the tabular case), the minimizer of $\Lcal_{\mathrm{DPO}}$ (when it exists, Proposition~1) is the policy of the two-stage pipeline: Bradley--Terry maximum likelihood $\hat r$, then $\pi^*_{\hat r}$ with the same $\beta$ and $\Lpiref$. \emph{Proof.} $\pi\mapsto\beta\log(\pi/\Lpiref)$ is a bijection from full-support distributions onto rewards modulo constants, with inverse $\hat r\mapsto\pi^*_{\hat r}$ (Proposition~2), and $\Lcal_{\mathrm{DPO}}(\pi)=\Lcal_{\mathrm{rew}}(\beta\log(\pi/\Lpiref))$.~\qedsymbol

Numerically (five responses, $6000$ comparisons, $\beta=0.5$) the two policies agree to $2.5\times10^{-8}$. With a neural policy the procedures differ: DPO optimizes within the class and never builds $\hat r$.

\textbf{4.\ Which objective is learned when.} \emph{Stage 1} (reward model): minimize $\Lcal_{\mathrm{rew}}(\varphi)$ over the comparisons; the unknown is the reward $r_\varphi$. \emph{Stage 2} (policy): with $\hat r$ frozen, maximize $J(\pi)=\E_\pi[\hat r]-\beta\KL(\pi\|\Lpiref)$; the unknown is the policy. \emph{DPO} does both at once. This is the content of the speaker's clarification of slide~25 (Section~\ref{sec:L6-questions}): the sigmoid loss belongs to stage~1, the objective $J$ to stage~2.
\end{beyond}

\begin{figure}[tb]
\centering
\begin{tikzpicture}
\begin{groupplot}[group style={group size=2 by 1, horizontal sep=14mm}, width=0.50\linewidth, height=0.40\linewidth,
    axis lines=left, axis line style={-{Latex[length=2mm,width=1.4mm]}, ink!70, line width=0.5pt},
    tick style={ink!70, line width=0.4pt}, every tick label/.append style={font=\sffamily\scriptsize, text=ink!85},
    label style={font=\sffamily\small, text=ink}, grid=major, major grid style={rule!70, line width=0.3pt},
    every axis plot/.append style={line width=1.3pt}, unbounded coords=jump, table/col sep=space,
    title style={font=\sffamily\small\bfseries, text=ink, yshift=-1mm}]
\nextgroupplot[title={(a) the KL-regularized optimum}, xmode=log, xmin=0.03, xmax=250, ymin=0, ymax=1, xlabel={$\beta$ (weight of the KL penalty)}, ylabel={probability $\pi^*_\beta(y_i)$}]
  \addplot[accent] table[x=beta, y=p1] {figures/data/L6_kl_policy.dat};
  \addplot[teal] table[x=beta, y=p2] {figures/data/L6_kl_policy.dat};
  \addplot[moss] table[x=beta, y=p3] {figures/data/L6_kl_policy.dat};
  \addplot[amber] table[x=beta, y=p4] {figures/data/L6_kl_policy.dat};
  \addplot[crimson] table[x=beta, y=p5] {figures/data/L6_kl_policy.dat};
  \node[font=\sffamily\fontsize{6}{7}\selectfont, text=accent, anchor=west] at (axis cs:38,0.45) {$y_1$};
  \node[font=\sffamily\fontsize{6}{7}\selectfont, text=teal, anchor=west] at (axis cs:38,0.30) {$y_2$};
  \node[font=\sffamily\fontsize{6}{7}\selectfont, text=moss, anchor=west] at (axis cs:38,0.15) {$y_3$};
  \node[font=\sffamily\fontsize{6}{7}\selectfont, text=amber, anchor=west] at (axis cs:38,0.075) {$y_4$};
  \node[font=\sffamily\fontsize{6}{7}\selectfont, text=crimson, anchor=west] at (axis cs:38,0.015) {$y_5$};
\nextgroupplot[title={(b) learning $r$ from comparisons}, xmode=log, ymode=log, xmin=80, xmax=70000, ymin=0.005, ymax=0.5, xlabel={number of comparisons $m$}, ylabel={RMS error of $\hat r$},
   legend style={font=\sffamily\fontsize{6.2}{7}\selectfont, at={(0.02,0.04)}, anchor=south west, draw=none, fill=white, fill opacity=0.85, text opacity=1}]
  \addplot[crimson, only marks, mark=*, mark size=1.8pt] table[x=n, y=rmse] {figures/data/L6_bt_consistency.dat};  \addlegendentry{simulation (60 runs)}
  \addplot[ink, dashed] table[x=n, y=asym] {figures/data/L6_bt_consistency.dat};  \addlegendentry{$\sqrt{\operatorname{tr}(I^{+})/(Km)}$}
\end{groupplot}
\end{tikzpicture}
\caption{\textbf{(a)} The optimum $\pi^*_\beta\propto\Lpiref e^{r/\beta}$ on five responses with rewards $r=(0,0.5,1,2,3)$ and reference $\Lpiref=(0.45,0.30,0.15,0.07,0.03)$: for large $\beta$ the model stays at the reference (right end), for small $\beta$ it puts all its mass on the best response. \textbf{(b)} Fitting a Bradley--Terry reward to simulated comparisons of $K=8$ responses: the error (up to a constant) falls like $m^{-1/2}$ and follows the Fisher-information prediction (dashed). Computed by \kbd{figures/gen\_data\_L6.py}.}
\label{fig:L6-bt}
\end{figure}

%% ------------------------------------------------------------------------------------
\section{The safety side: from human feedback to documents}\label{sec:L6-safety}
\tsm{53:35}
\subsection{What is wrong with RLHF}\label{sec:L6-rlhflimits}\tsm{53:35}
The speaker is going more slowly than planned; what he wants to discuss is the safety aspect of the training. One issue is that RLHF is \emph{not scalable}: the number of human answers used ranges from $100{,}000$ to a million, which is not enough, and more is needed as the models grow. The other is that models with too much RLHF do not necessarily show the best behavior: they optimize the immediate reward of a human and so can become \term{sycophantic}, telling you that you are an amazing person and that all your ideas are genius. We do not want to optimize for the split-second reaction of a human; we would like to optimize for longer-term things and to think about consequences. So the way various companies approach this is to formulate general ideas, a document that summarizes what we want these models to do and how we want them to behave. The two famous examples are \emph{Claude's Constitution} and OpenAI's \emph{Model Spec}, and the difference between them is, he says, philosophically interesting.

\begin{beyond}[Scale of the preference data, sycophancy, over-optimization]\label{box:L6-sycophancy}
The numbers are of the right order: InstructGPT's reward model used $33$k prompts with $K=4$ to $9$ ranked completions each, so $\binom K2=6$ to $36$ comparisons per prompt and between about $2\times10^5$ and $1.2\times10^6$ comparisons \citep{ouyang2022}; Llama~2 reports ``over 1 million binary comparisons'' \citep{touvron2023}. Sycophancy is documented in five state-of-the-art assistants; in existing human preference data, responses agreeing with the user's views are more likely to be preferred, and optimizing against preference models can sacrifice truthfulness \citep{sharma2023}. Over-optimization of an imperfect reward is in \citet{gao2023} (see the box \emph{Bradley--Terry, the KL leash and DPO}); preference-based reward learning for deep RL goes back to \citet{christiano2017}.
\end{beyond}

\subsection{Two documents}\label{sec:L6-docs}\tsm{54:35}
Claude's Constitution is very philosophical. Its main author is Amanda Askell, a philosopher with a PhD from NYU (the speaker first calls her one of the authors, then corrects himself). The document is $84$ pages long and talks a lot about Claude's \emph{character}: it focuses not on what Claude should or should not do but on what kind of character, what persona, it should have, and spends much time describing it. Claude should act like a best friend of a human, someone who is an expert engineer and scientist, perhaps also a lawyer and a doctor, and also a good friend: how would such a person act? OpenAI's approach is, by comparison, more \emph{behavior}-focused. Anthropic's claim is that they find character training more effective for alignment: when they add more examples of models thinking carefully through moral and ethical dilemmas before deciding, in mid-training, in supervised fine-tuning and in the RL stage, the model ends up more aligned in a more diverse set of situations, also out of distribution.

Rather like the laws of robotics of the famous science-fiction author, the document has a \emph{hierarchy} of rules (Figure~\ref{fig:L6-ladders}, left). The rules about oversight, about Claude's being steerable and inspectable by humans, are the most important. Next in priority is that Claude be ethical; next, that it adhere to what Anthropic thinks it should do; after that it should be helpful, to what the user asks or, through the API, to what a third-party company wants, and after that to what the user wants. The speaker finds this structure interesting; the document's point of view, he says, is that it is more important for Claude to develop general ethical principles than just to listen to whatever Anthropic tells it. The constitution is organized into sections and he shows some excerpts.

\subsection{The Model Spec and deliberative alignment}\label{sec:L6-spec}\tsm{59:05}
\term{Deliberative alignment} is OpenAI's approach, and again there is a hierarchy: \emph{root, system, developer, user, guideline} (Figure~\ref{fig:L6-ladders}, right). \emph{Root} holds the most important rules, the principles under which the model should operate: for example, that it should not seek power and should not try to deceive humans. Then come what OpenAI wants (\emph{system}), what the third-party provider wants (\emph{developer}), what the actual user wants, and finally the tips for being helpful, reasonable and having common sense (\emph{guidelines}). ``That is how I understand these levels.''

\begin{figure}[tb]
\centering
\begin{tikzpicture}[x=1cm, y=1cm, >={Latex[length=2mm,width=1.4mm]},
    rung/.style={card=#1, text width=5.3cm, align=left, inner sep=4pt},
    lab/.style={font=\sffamily\fontsize{6.6}{7.4}\selectfont, text=slate, align=left}]
  \node[font=\sffamily\small\bfseries, text=ink] at (3.0,6.45) {Claude's Constitution};
  \node[rung=crimson, minimum height=1.0cm] (c1) at (3.0,5.5) {\textbf{1\ Broadly safe}: not undermining appropriate human oversight of AI};
  \node[rung=plum, minimum height=1.0cm] (c2) at (3.0,4.15) {\textbf{2\ Broadly ethical}: good values, honesty, avoiding harm};
  \node[rung=amber, minimum height=1.0cm] (c3) at (3.0,2.8) {\textbf{3\ Compliant with Anthropic's guidelines}};
  \node[rung=moss, minimum height=1.0cm] (c4) at (3.0,1.45) {\textbf{4\ Genuinely helpful}: to operators and users};
  \foreach \a/\b in {c1/c2, c2/c3, c3/c4} \draw[->, slate, line width=0.8pt] (\a.south) -- (\b.north);
  \node[lab, text width=5.6cm] at (3.0,0.4) {``in the order in which they are listed'', but ``holistic rather than strict''; hard constraints apply separately};
  \node[font=\sffamily\small\bfseries, text=ink] at (9.8,6.45) {OpenAI Model Spec: chain of command};
  \node[rung=crimson, minimum height=0.8cm] (m1) at (9.8,5.6) {\textbf{Root}: fundamental rules that nothing below can override};
  \node[rung=plum, minimum height=0.8cm] (m2) at (9.8,4.55) {\textbf{System}: what OpenAI wants};
  \node[rung=amber, minimum height=0.8cm] (m3) at (9.8,3.5) {\textbf{Developer}: what the third-party provider wants};
  \node[rung=moss, minimum height=0.8cm] (m4) at (9.8,2.45) {\textbf{User}: what the actual user wants};
  \node[rung=slate, minimum height=0.8cm] (m5) at (9.8,1.4) {\textbf{Guideline}: tips for being helpful; overridable};
  \foreach \a/\b in {m1/m2, m2/m3, m3/m4, m4/m5} \draw[->, slate, line width=0.8pt] (\a.south) -- (\b.north);
  \node[lab, text width=5.6cm] at (9.8,0.4) {the lecture's examples of root rules: do not seek power, do not deceive humans (not checked)};
\end{tikzpicture}
\caption{The two ``ladders''; higher rungs prevail in a conflict. \textbf{Left:} the four core properties of Claude's Constitution, in the order and with the caveats of the document (see the box \emph{What the two documents say, as far as I could check}). \textbf{Right:} the chain of command of the Model Spec, in the order recorded in OpenAI's changelog of 12~September 2025, ``root'' as defined there; the glosses of the other levels are the lecture's and that of guidelines is from search excerpts (I could not open the Model Spec itself).}
\label{fig:L6-ladders}
\end{figure}

\begin{beyond}[What the two documents say, as far as I could check]\label{box:L6-docs}
\emph{The Constitution} \citep{anthropic2026constitution}, published on 21~January 2026 under a CC0 licence: I read the PDF, $84$ pages and about $30{,}000$ words; authors Amanda Askell and Joe Carlsmith (lead), Chris Olah, Jared Kaplan, Holden Karnofsky, ``several Claude models, and many other contributors''. It lists \emph{broadly safe}, \emph{broadly ethical}, \emph{compliant with Anthropic's guidelines}, \emph{genuinely helpful} and says: ``In cases of apparent conflict, Claude should generally prioritize these properties in the order in which they are listed'', the prioritization being ``holistic rather than strict''. Two refinements of the summary: oversight is not obedience (``being overseeable \dots does not mean blind obedience, including towards Anthropic''); and ``operators before users'' is not the fourth rung but a separate principal hierarchy (Anthropic, operators, users) with trust ``in roughly the order given above \dots This is not a strict hierarchy''. The document says it ``plays a crucial role in our training process'', and explains rather than orders because ``if we train Claude to exhibit even quite narrow behavior, this often has broad effects on the model's understanding of who Claude is''.

\emph{The Model Spec} \citep{openai2026spec}: the order \emph{Root, System, Developer, User, Guideline} is that of the changelog entry of 12~September 2025, which renamed the top level from Platform to Root and placed it above System, seen through a changelog aggregator; my tools could not fetch the Model Spec pages, so the definitions of the levels and the lecture's examples of root rules are unchecked. The hierarchy belongs to the Model Spec; \emph{deliberative alignment} \citep{guan2024} is the training method that uses it. As OpenAI describes it, a model trained only to be helpful generates (prompt, chain of thought, answer) triples with the relevant part of the specification in its system prompt; a judge that also sees the specification filters them; the model is fine-tuned on them \emph{with the specification removed from the prompt}; reinforcement learning then follows, with a reward from a judge that has access to the specification. This matches the lecture.
\end{beyond}

\subsection{Questions}\label{sec:L6-questions}\tsm{1:00:30}
\paragraph{A question about slide 25.}\tsm{1:00:30}
\begin{discussion}[What is the reward in the loss?]
A student: on slide~25 the \emph{loss of the reward} is written, but the reward is not defined. The speaker: it is the expectation of the logarithm of $\sig$ of the difference of the rewards of the two answers. The student: the definition of $R$ itself, $R_\varphi$. The speaker: you do not know how to find $R$; you want to deduce it by inverting this relationship, to \emph{learn} it; that is part of the difficulty, and there are various algorithms, among them DPO, which sidesteps computing the reward. The student: maybe the reward is also in the loss. The speaker: ``So what did I write here? It's maybe a little bit confusing \ldots{} You're right \ldots{} I'm a little bit confused by why I put this here.'' What he wanted to write is that one wants to \emph{maximize} the reward minus the penalty (the objective $J$ of the slide). And separately: ``you want to learn the reward. How do you learn the reward? By minimizing this loss''; that is the procedure for learning the reward, the unknown; ``now once you have learned the reward you can optimize, by maximizing $J$''.
\end{discussion}
\noindent So the slide mixed the objectives of two stages; the box \emph{Bradley--Terry, the KL leash and DPO} (item~4) states which is learned when, and that DPO does both at once.

\paragraph{Are the documents in the system prompt?}\tsm{1:03:30}
\begin{discussion}[How the Constitution and the Model Spec are used]
Not through the system prompt; they are used in a few other ways. For the Constitution: (i)~in \textbf{mid-training}, on a much larger corpus than the Constitution itself, which elaborates it with many examples of models meeting dilemmas, thinking through parts of the constitution and deciding, trained by next-token prediction; (ii)~in \textbf{SFT} on a lot of that; (iii)~in \textbf{RLAIF}: a different model that carefully reads the constitution is given two answers to a prompt and chooses the one that adheres more to it; it is an AI judge, but with specific instructions on how good a response is from the point of view of the constitution, and the model is trained by RL against it. The Model Spec approach is similar: (1)~generate prompts that touch important aspects of the spec; (2)~give the model the spec, so that it reasons about the problem \emph{while thinking about the spec}; (3)~have the answer judged by an AI judge that also reasons about the spec; (4)~\emph{remove} the spec and train on these examples, so that the model adheres to the spec even when it is not thinking about it.
\end{discussion}

\paragraph{Is there evidence that it works?}\tsm{1:07:10}
A student asks whether both approaches are just companies improvising, and whether there is evidence, experimental or theoretical, that this works to some measure. The speaker: there are papers showing that, on various examples of dilemmas, after this type of training the models tend to give better answers. The deliberative-alignment paper has a bunch of experiments; and in \emph{Teaching Claude Why} (May 2026) they show that a Claude trained just on good answers does not generalize well to other situations, while one trained also on reasoning about the ethical dilemmas involved and the relevant parts of the constitution gives good answers even to questions not very similar to those seen in training. How much, if at all, this generalizes to the dilemmas that arise with larger capabilities is hard to say: ``a very much experimental kind of thing''.

\begin{beyond}[What ``Teaching Claude Why'' reports]\label{box:L6-why}
I read Anthropic's post \citep{kutasov2026} (8~May 2026) through page summaries and search excerpts. It concerns \emph{agentic misalignment} in fictional scenarios, such as a model blackmailing an engineer to avoid shutdown. Training on about $10$k transcripts ($30$M tokens) filtered only to remove responses that took the bait lowered the misaligned behavior from $22\%$ to $15\%$; when the responses were rewritten to contain explicit reasoning about values, to about $3\%$. Chats in which the assistant advises a \emph{user} facing an ethical dilemma matched this with only about $3$M tokens and generalized better; documents about the constitution plus fictional stories about admirable AIs, unrelated to the evaluation, cut blackmail from $65\%$ to $19\%$ in one configuration. The authors stress that training on the evaluation distribution lowered the measured behavior but \emph{not} held-out audit metrics, that experiments were mostly on Sonnet- and Haiku-class models, and that a model may memorize the constitution without internalizing it. So the lecture's summary has the right direction, but what was measured is behavior in agentic-misalignment evaluations, not ``ethical dilemmas'' in general. The earlier AI-feedback method with a list of principles is \citet{bai2022}.
\end{beyond}

\paragraph{How was it baked in? Announcement.}\tsm{1:09:15}
A last question: how exactly were the Constitution and the ethics baked into the model's identity? Again through mid-training on many examples of thinking generated by LLMs from the Constitution, and through the later RL from AI feedback, in which the AI judges whether a response adheres to the Constitution. There is some human input, but it is fairly small compared with the amount of AI-generated material. The speaker closes by saying that many things happen on Friday, so he is not sure that he can hold office hours that day, and that he will make another announcement.\tsm{1:10:30}

%% ------------------------------------------------------------------------------------
\tocsection{Take-aways and exercises}
\begin{openq}[Questions this lecture leaves open]
\begin{itemize}
  \item Which monitoring layers are deployed for \emph{agentic} behavior, and how long do chain-of-thought monitors stay informative when the monitored model is trained with them in the loop?
  \item Teaching \emph{reasons} seems to generalize better than teaching behaviors. Is there a model in which a short principle plus reasoning beats a long list of demonstrations, in the spirit of Lecture~\ref{ch:L3}? What is lost when a specification is removed from the context (deliberative alignment as context distillation)?
  \item Can the expressivity results for chains of thought be turned into statements about \emph{learning}: how many examples teach a transformer to use a scratchpad? How does one tell that a constitution is internalized rather than memorized?
\end{itemize}
\end{openq}

\begin{takeaways}
\begin{itemize}
  \item An LLM is a chain of vector spaces; each transformer layer adds \emph{causal attention} to an MLP, at one width. The last vector gives scores $z_v=\ip{u_v}{h}+b_v$, probabilities $\softmax(z/\tau)$ and a sampled token, appended to the text; the loop repeats, and a cache avoids recomputation.
  \item The model sees a hidden system prompt, the history and a hidden chain of thought (a hard-coded think token, a learned end token), which buys serial computation. An agent is an LLM plus a harness with its state in files; four monitoring layers exist or are planned.
  \item Pre-training minimizes $-\log P_\theta(x_t\mid x_{<t})$; mid-training selects better data; post-training is SFT, RLHF/RLAIF and RL with verifiable rewards. The compute share of RL rose from $\approx0.2\%$ (DeepSeek-V3) to above $10\%$ of the pre-training cost (V3.2); text is nearly exhausted, and $20$ tokens per parameter is, in the speaker's words, probably starting to break at the frontier.
  \item The next-token loss is code length: good prediction is compression, finding short reasons, not shallow statistics; but it is not what we want the assistant to do.
  \item $\Pr(w\succ l)=\sig(r(w)-r(l))$: the reward is known up to a constant and learned by a logistic loss; the policy maximizes $\E[r]-\beta\KL(\pi\|\Lpiref)$, solved by $\pi^*\propto\Lpiref e^{r/\beta}$; DPO substitutes this into the likelihood and skips the reward.
  \item RLHF does not scale and breeds sycophancy; Claude's Constitution and OpenAI's Model Spec (with deliberative alignment: reason about the spec, then drop it) are trained in by mid-training, SFT and RLAIF; the evidence is experimental.
\end{itemize}
\end{takeaways}

\begin{exercise}[Problems]
\begin{enumerate}
  \item \emph{Temperature.} For $z=(3,1,0)$ compute $p_\tau$ and its entropy (nats) at $\tau=0.5,1,2$. Prove $\dd H/\dd\beta=-\beta\Var_\beta(z)$ for $\beta=1/\tau$, and the limits $\tau\to0^+$, $\tau\to\infty$.
  \item \emph{Logit versus angle.} For $h=(1,0)$, $b=0$, $u_1=(0.5,0.2)$, $u_2=(2,2)$ find the angles and the probabilities at $\tau=1$. Show that for equal norms the order of the probabilities is that of the cosines, and that $\norm h$ acts as an inverse temperature.
  \item \emph{Parameters, compute, cache.} Count the parameters of a dense transformer with $L$ layers, width $d$, MLP width $4d$ and a tied embedding of $V$ tokens (no biases). For $L=96$, $d=12288$, $V=50257$ and $2048$ positions compare with $175$ billion, and compute $6ND$ for $D=3\times10^{11}$. How many bytes per token does a key--value cache need in 16-bit numbers, and for $10^5$ tokens? How many positions are processed with and without the cache when $T=500$ tokens follow a prompt of $P=1000$ (the last generated token is not fed back)?
  \item \emph{Bradley--Terry.} (a) Show $\Lcal_{\mathrm{rew}}(r+c)=\Lcal_{\mathrm{rew}}(r)$. (b) Answers $A,B,C$ with the comparisons $A\succ B$, $B\succ C$, $A\succ C$ ($3$ times each) and no others: show that no maximum-likelihood reward exists. (c) Add one comparison $C\succ A$ and find the maximizer (with $\sum r=0$) numerically.
  \item \emph{The leash.} Prove $J(\pi)=\beta\log Z-\beta\KL(\pi\|\pi^*)$. For $\Lpiref=(0.5,0.3,0.2)$, $r=(0,1,2)$, $\beta=1$ compute $\pi^*$, $\E_{\pi^*}r$, $\KL(\pi^*\|\Lpiref)$ and $\max J$; repeat with $\beta=0.5$. What happens as $\beta\to0$ and $\beta\to\infty$?
  \item \emph{DPO.} Derive the DPO loss from the Bradley--Terry model and the optimum of Exercise~5. For $\beta=0.1$ and log-ratios $\log(\pi_\theta/\Lpiref)$ equal to $1.2$ (winner) and $-0.8$ (loser), compute the loss.
  \item \emph{Chinchilla arithmetic.} With $N=D/20$ and $C=6ND$ find $N$ and $C$ for $D=1.5\times10^{13}$ and $D=3.2\times10^{14}$ tokens. DeepSeek-V3 has $37$B active parameters and $14.8$T tokens: how many tokens per active parameter, and what is $6ND$?
  \item \emph{Scratchpads.} The product of $n$ permutations of $\{1,\dots,5\}$ can be computed by a scratchpad of $n$ steps, each using only the previous token and the next input. Show that associativity also allows a balanced tree of depth $\lceil\log_2n\rceil$, and verify with $n=64$. Which of the two matches a fixed-depth transformer's single pass, and what does a chain of thought trade for what?
  \item \emph{Prediction and compression.} A model has loss $2.0$ nats per token and a token averages $4$ bytes. Find its perplexity, its bits per byte and its compression ratio, and compare with $8.3\%$ for Chinchilla~70B on enwik9. A Bayes mixture of $7$ predictors with a uniform prior is at most how many bits worse than the best one?
\end{enumerate}
\end{exercise}


